\section{Conclusion}

The Critical Line Algorithm computes the whole mean--variance frontier by walking its
piecewise-linear path from the maximum-return to the minimum-variance portfolio.
\pkg{cvxcla} takes each step with one block-eliminated KKT solve, resolves ties by a
least-index rule under which the walk provably terminates (Proposition~\ref{prop:termination}),
and reaches the covariance only through a five-member operator protocol. A factor risk
model therefore yields the same turning points as the dense matrix it represents; with $2560$ assets and ten factors, on a market where every asset is free at some point, it does so about $30$ times faster, in under $1$~MiB of working memory beyond the stored frontier. There a general interior-point solver given the same factor form returns a grid of $100$ frontier portfolios at a similar cost, but only to its solver tolerance, with variances off in the third digit; on the S\&P~500, where few assets are free at once, the grid takes more than twenty times as long as the trace. The trace returns every turning point, exact up to round-off. The gain is one of scale and structure. Below about $200$
assets the per-step overhead of \proglang{Python} dominates, and \pkg{qpOASES}, which
follows the same path, is faster.
